\section{Paired walks, sorted layers, and shared axes}\label{sec:layers}
\subsection{Paired decomposition and the exit convention}\label{sec:paired-walks}
At each nonfixed vertex choose a matching from incoming to outgoing
edge copies and force its reversed matching at the mate vertex. At a
fixed vertex, pair outgoing copies by a fixed-point-free involution;
identifying incoming copies with their reversals gives the transition
matching. This is possible by Section~\ref{sec:base}. The resulting cycles
come in distinct reversed pairs. For if reversal fixed one cyclic
edge order, its reflection would fix an edge or exchange two consecutive
edges, yielding a fixed transition. Both are excluded. Fixed tie orders
make this construction deterministic. It is an elementary adaptation
of paired-decomposition principles, not a separate contribution.

Lift each closed walk to windows $[x,e)$ in the periodic text formed
by its cyclic up letters. Up steps increase $e$ and down steps increase
$x$; a turn with $P$ up steps advances by $(P,P)$. A word occurrence
is recorded if a visited window contains it, counted once per starting
residue modulo $P$. The telescoping up-step argument of
\src{\texttt{lem:recorded}} gives total recorded count $m(s)$,
even when $|s|>P$.

Group texts up to translation, reduce to least period $p$, and sort
primitive passages by normalized exits. More explicitly, if original
walk $j$ has turn $d_jp$ and exit function $Z_j$, sort the multiset
\[
\{Z_j(x+rp)-rp: j,\ 0\le r<d_j\}
\]
at each $x$. Its order statistics are layer exits $z_i(x)$.
We index higher layers first. Put
\begin{equation}\label{eq:profiles}
f_i(x)=z_i(x-1),\quad l_i(x)=z_i(x),\quad
z_i(x-1)\le z_i(x),\quad z_i(x-1)\ge x,\quad
z_i(x+p)=z_i(x)+p.
\end{equation}
A layer visits exactly $[x,e)$ for $f_i(x)\le e\le l_i(x)$, then
takes a down step. The source sorting proof preserves the phase-labelled
edge multiset: down steps are determined by exits, and vertical-edge
multiplicities by entry-minus-exit cumulative counts. Each layer has
$p$ up steps and receives an additional budget $p$; the budgets sum
to $W$. Every layer window spells a word of $V$.

The convention for occurrences is crucial. At a specified start $t$,
a layer records $[t,t+h)$ exactly when its \emph{exit}
$z_i(t)\ge t+h$. Thus sorting preserves phase-specific recorded counts,
not merely aggregate counts of word labels. Its entries decide whether
that recorded word is also an actual vertex at start $t$.

\subsection{Text orbits and adjacent pairs}\label{sec:axes}
A nonself text orbit comprises two translation classes exchanged by
$\rho$. Their sorted layers pair rank by rank. Choose the canonical
text to be the greatest forward rotation among \emph{both} orientations.
A basic unit consists of one representative layer and its reflected
mate, with budget $2p$.

For a self-reversed primitive text fix
\begin{equation}\label{eq:reflection}
\mathcal R(x,e)=(c-e,c-x),\qquad
\cm(A(j))=A(c-1-j).
\end{equation}
There are an even number of primitive passages, inherited from distinct
reversed walk pairs. Reflection preserves the order of staircase
profiles: it sends the upper region of a path to the upper region of
the reflected path. The ordered decomposition is uniquely determined
by its exit multisets. Its phase-labelled multiset is invariant under
$\mathcal R$, so every sorted profile is individually fixed under this
reflection, modulo period translation. This does \emph{not} imply even
multiplicity of each profile.

\begin{lemma}[Shared-axis adjacent pairs]\label{lem:axes}
Pair the sorted layers of a self text from the top in adjacent pairs.
Each pair shares both reflection-axis vertices in a primitive turn,
and is therefore already connected. The two profiles may be different.
Removing full pairs preserves all fixed-phase even tails.
\end{lemma}
\begin{proof}
For a nonempty interval $[t,t+h)$ with $2t+h\equiv c\pmod p$,
reflection fixes its alignment. It pairs its recorded occurrences on
the original distinct reversed walks without fixed occurrences.
Consequently their number is even. The exit convention makes this
number exactly
\begin{equation}\label{eq:even-tail}
\#\{i:z_i(t)\ge t+h\},
\end{equation}
which sorting preserves. Since exits are ordered, every even upper
tail consists of whole top-aligned adjacent pairs. Whenever the upper
member visits this axis vertex, the lower member records it too; its
entry is no higher than the upper entry, so it actually visits the same
vertex. The two threshold indicators agree for every pair, and removing
a full pair preserves evenness.

Along one primitive turn, $x+e$ increases by one at every step and by
$2p$ overall. It therefore reaches both classes $c$ and $c+p$ modulo
$2p$, at actual vertices. These are the two fixed-axis vertices, where
the preceding argument applies. If the vertex is empty, its upper
entry is $x$, and the lower entry, lying between $x$ and that entry,
is also $x$; both visit $\eps$. No occurrence count of the empty word
is needed. When $p=1$, every window is an axis window and the turn has
exactly two window visits; the same argument includes this case.
\end{proof}

\subsection{Spliced paths for hosting}\label{sec:splice}
Let $D_1,D_2$ be a self pair and let $U,V$ be its two shared axes in
cyclic order. Form $E$ from $D_1[U,V]$ followed by
$D_2[V,U+(p,p)]$, and form $F$ from the complementary arcs.
These are monotone period-$p$ paths with $p$ up steps each; their
labelled union equals that of $D_1,D_2$. Reflection exchanges the
half-period arcs, so $\mathcal R(E)=F$ at the level of windows and
phase-labelled edges. Coincident profiles use distinct labels.
This is not an assertion about the original base edge-copy involution
after sorting. Later rootedness uses shared vertices, and mirrored
added walks receive new paired edge copies. Keep the original sorted
profiles for all band geometry; use $E,F$ only as host paths.

The exchange-of-halves operation has the antecedent in
\cite[Section 5]{gk}. That operation assumes the two meetings.
Lemma~\ref{lem:axes} is what establishes them for the ordered profiles here.
